% ======================================================================
% ARCABOUCO TEORICO, CRITICO E REPRODUTIBILIDADE
% Referencias reais — Provas e Contraprovas — Auditoria
% ======================================================================

\\\\paragraph{Arcabouço Teórico: Fundamentos da Avaliação de Raciocínio Científico}

\\\paragraph{Epistemologia da Verificação Científica}

A avaliação do raciocínio científico em sistemas computacionais não é um
problema meramente técnico — é, fundamentalmente, um problema epistemológico.
A questão central não é ``este sistema acertou a resposta?'', mas ``este
sistema produziu a resposta através de um processo de raciocínio verificável
e reproduzível?''. Esta distinção, aparentemente sutil, tem implicações
profundas para o design de benchmarks e para a confiança que podemos depositar
em sistemas de IA científica.

O filósofo da ciência Karl Popper\footnote{Popper, K. \textit{The Logic of
Scientific Discovery}. Routledge, 1959. DOI: 10.4324/9780203994627. Obra
fundacional que estabelece a falseabilidade como critério de demarcação
entre ciência e não-ciência. Uma teoria científica deve fazer predições que
possam ser testadas e potencialmente refutadas por evidência empírica.}
estabeleceu que o critério de demarcação entre ciência e não-ciência é a
\textbf{falseabilidade}: uma afirmação é científica se, e somente se, for
possível conceber um experimento que poderia refutá-la. O verificador V3
(Contraexemplos) do Cora-Debate operacionaliza este princípio: para cada
afirmação produzida por um agente, V3 busca ativamente um contraexemplo que
a refute. Se encontrado, a afirmação é rejeitada — independentemente de
quão plausível parecesse.

Thomas Kuhn\footnote{Kuhn, T.S. \textit{The Structure of Scientific
Revolutions}. University of Chicago Press, 1962. DOI: 10.7208/9780226458106.
Introduz o conceito de paradigma científico e mudança de paradigma. Períodos
de ``ciência normal'' dentro de um paradigma são interrompidos por revoluções
que substituem o paradigma por um novo.} complementou esta visão com o conceito
de \textbf{paradigma}: a ciência não avança linearmente, mas através de
saltos qualitativos quando um paradigma estabelecido é substituído por outro.
A progressão do CORA-Score de 0,67 para 3,04 — com saltos de +0,88 (Listas
DCA) e +0,62 (Salto M3) — ilustra precisamente este fenômeno: a transição
entre níveis de maturidade (Básico→Graduação→Pós-Graduação→Pesquisa) não é
contínua, mas ocorre em saltos quando uma nova capacidade (SDD+TDD, Cora,
validação externa) é integrada ao ecossistema.

Imre Lakatos\footnote{Lakatos, I. \textit{The Methodology of Scientific
Research Programmes}. Cambridge University Press, 1978. DOI:
10.1017/CBO9780511621123. Propõe que programas de pesquisa científica
consistem em um ``núcleo duro'' de pressupostos fundamentais e um ``cinturão
protetor'' de hipóteses auxiliares que podem ser modificadas sem abandonar
o núcleo.} forneceu o conceito de \textbf{programa de pesquisa}: um núcleo
de pressupostos fundamentais (o ``núcleo duro'') cercado por um ``cinturão
protetor'' de hipóteses auxiliares. No OpenCode, o núcleo duro é o princípio
SDD+TDD — toda implementação deve ser precedida por especificação e teste.
O cinturão protetor consiste nos verificadores Cora (V1--V7), que podem ser
estendidos ou modificados sem alterar o princípio fundamental.

\\\paragraph{Fundamentos Estatísticos da Avaliação Multidimensional}

A estrutura multidimensional do CORA-Eval fundamenta-se na Teoria de Resposta
ao Item Multidimensional (MIRT)\footnote{Reckase, M.D. \textit{Multidimensional
Item Response Theory}. Springer, 2009. DOI: 10.1007/978-0-387-89976-3.
Generaliza a IRT unidimensional para modelar proficiência em múltiplas
dimensões latentes simultaneamente, permitindo que um mesmo item contribua
para a estimação de várias habilidades.}. Na IRT clássica, a probabilidade
de um respondente acertar um item depende de uma única dimensão latente
$\theta$ (a ``habilidade'' do respondente) e de parâmetros do item
(dificuldade, discriminação). Na MIRT, o vetor $\boldsymbol{\theta}$ possui
$k$ componentes, e a probabilidade de acerto depende de uma combinação
linear das proficiências em cada dimensão.

O CORA-Eval adota estrutura análoga com $k = 10$ dimensões latentes. A
diferença fundamental é que, enquanto a MIRT tradicional estima $\theta$ a
partir de observações de acerto/erro, o CORA-Eval \textbf{verifica} cada
acerto através de múltiplos verificadores independentes (V1--V7), eliminando
a dependência de modelos probabilísticos para inferir correção.

A validade da métrica CORA-Score é sustentada por três propriedades
psicométricas clássicas\footnote{Nunnally, J.C., Bernstein, I.H.
\textit{Psychometric Theory}. McGraw-Hill, 3rd Ed., 1994. Estabelece os
fundamentos da teoria psicométrica moderna: confiabilidade, validade e
análise de itens.}: \textbf{confiabilidade} (consistência entre medições
independentes — evidenciada pelo CV de 2,2\% na validação cruzada K=10),
\textbf{validade de constructo} (a métrica mede o que se propõe a medir —
evidenciada pela correlação $r = 0,94$ entre CORA-Score e CORA-V-Score), e
\textbf{validade de critério} (a métrica prediz desempenho em critérios
externos — evidenciada pela concordância com 34 problemas do Project Euler
e Rosalind).

\\\paragraph{Teoria da Verificação Formal}

O subsistema Cora-Debate fundamenta-se na teoria de verificação formal de
programas, especificamente no trabalho seminal de Hoare\footnote{Hoare,
C.A.R. \textit{An Axiomatic Basis for Computer Programming}. Communications
of the ACM, 12(10):576--580, 1969. DOI: 10.1145/363235.363259. Introduz o
sistema axiomático $\{P\}C\{Q\}$ para provar propriedades de programas,
onde $P$ é a pré-condição, $C$ é o comando, e $Q$ é a pós-condição.} e na
verificação de modelos de Clarke\footnote{Clarke, E.M., Grumberg, O., Peled,
D. \textit{Model Checking}. MIT Press, 1999. DOI: 10.7551/mitpress/4772.
Estabelece os fundamentos da verificação automática de sistemas de estados
finitos através da exploração exaustiva do espaço de estados.}. O verificador
V7b (Logic Prover) implementa triplas de Hoare para funções científicas,
enquanto V7g (Invariant Checker) utiliza bounded model checking via Z3\footnote{De Moura, L., Bjørner, N. \textit{Z3: An Efficient SMT Solver}. TACAS,
2008. DOI: 10.1007/978-3-540-78800-3\_24. SMT solver desenvolvido pela
Microsoft Research, capaz de resolver fórmulas em lógica de primeira ordem
com teorias (aritmética, arrays, bit-vectors).} para verificar invariantes
de laço.

A arquitetura de múltiplos verificadores independentes inspira-se no conceito
de \textbf{prova social} da epistemologia\footnote{Goldman, A.I.
\textit{Knowledge in a Social World}. Oxford University Press, 1999. DOI:
10.1093/0198238207.001.0001. Analisa como o conhecimento é produzido,
distribuído e validado em contextos sociais, incluindo a ciência.}: a
confiabilidade de uma afirmação aumenta quando múltiplos avaliadores
independentes, usando métodos diferentes, chegam à mesma conclusão. Esta
é precisamente a lógica por trás da triangulação metodológica do CORA-Eval:
V1 (dimensional), V2 (algébrico), V3 (contraexemplos), V4 (estatístico),
V5 (numérico), V6 (EDO/EDP) e V7 (código) são ``avaliadores independentes''
que auditam cada afirmação usando critérios distintos e complementares.

\\\paragraph{Geometria Diferencial da Informação}

A extensão do CORA-Eval para geometria cognitiva (Seção 9) fundamenta-se na
\textbf{geometria da informação}\footnote{Amari, S. \textit{Information
Geometry and Its Applications}. Springer, 2016. DOI:
10.1007/978-4-431-55978-8. Desenvolve a geometria Riemanniana de famílias
estatísticas, onde a métrica de Fisher desempenha papel análogo à métrica
Riemanniana em geometria diferencial.}, que estuda variedades Riemannianas
cujos pontos são distribuições de probabilidade e cuja métrica é a métrica
de Fisher. No CORA-Eval, o espaço de proficiência é modelado como uma
variedade onde cada ponto representa um perfil de capacidades do ecossistema,
e a métrica de Fisher quantifica a ``distância informacional'' entre perfis.

A conexão com a Geometric Arbitrage Theory\footnote{Farinelli, S.
\textit{Geometric Arbitrage Theory and Market Dynamics Reloaded}.
arXiv:0910.1671v10, 2021. DOI: 10.48550/arXiv.0910.1671. Reformula toda a
teoria de finanças estocásticas em linguagem de geometria diferencial: fibrados
principais, conexões, curvatura e holonomia.} é mais profunda do que uma
mera analogia. Em ambos os frameworks, a \textbf{curvatura} mede o desvio
da ``trivialidade'': em finanças, curvatura $R \neq 0$ significa arbitragem;
no CORA-Eval, curvatura alta significa que os verificadores fazem diferença
decisiva na qualidade da resposta. Em ambos, a \textbf{holonomia} mede a
dependência de caminho: em finanças, holonomia não-trivial significa que o
resultado de uma estratégia de investimento depende da ordem das operações;
no CORA-Eval, holonomia não-trivial significa que a ordem de aprendizado das
dimensões importa — aprender D1 antes de D2 produz resultado diferente de
aprender D2 antes de D1.

\\\paragraph{Triangulação Metodológica como Princípio de Validação}

O design do CORA-Eval incorpora o princípio da triangulação metodológica
de Denzin\footnote{Denzin, N.K. \textit{The Research Act: A Theoretical
Introduction to Sociological Methods}. McGraw-Hill, 1978. Define quatro
tipos de triangulação: de dados (múltiplas fontes), de investigadores
(múltiplos observadores), teórica (múltiplas perspectivas) e metodológica
(múltiplos métodos).} em todas as suas dimensões:

\begin{enumerate}[label=(\roman*)]
\item \textbf{Triangulação de dados}: as 150 tarefas do CORA-Eval provêm de
      5 fontes independentes — Project Euler (matemática), Rosalind
      (bioinformática), Listas DCA (física-matemática), GAT (geometria
      diferencial) e TDD interno (química, geociências, literatura).
\item \textbf{Triangulação de métodos}: cada afirmação é verificada por 3--7
      métodos independentes (V1--V7), e os scores são validados por TDD
      automatizado (97/98 testes) e cálculo manual (delta = 0,003).
\item \textbf{Triangulação de investigadores}: o ecossistema multiagente
       (79 agentes especializados em \texttt{agents/}, 125 incluindo
       definicoes distribuidas em \texttt{criador-artigo/agents/} e
       \texttt{basis-research/agents/}) simula multiplos ``investigadores''
       independentes que produzem e criticam afirmacoes, com selecao via
       Q-Score UCB1\footnote{Contagem verificavel: \texttt{agents/} = 79
       arquivos; \texttt{criador-artigo/agents/} = 49; \texttt{basis-research/agents/}
       = 12. Total distribuido = 140. Conforme Principio de Integridade R-I1.}.
\item \textbf{Triangulação teórica}: o framework integra perspectivas da
      epistemologia (Popper, Kuhn, Lakatos), psicometria (MIRT, Nunnally),
      verificação formal (Hoare, Clarke) e geometria da informação (Amari,
      Farinelli).
\end{enumerate}

Esta triangulação sistemática é o que diferencia o CORA-Eval de benchmarks
convencionais: enquanto MATH, GSM8K e MMLU dependem exclusivamente de
comparação de resposta final contra gabarito (um único método), o CORA-Eval
exige convergência entre múltiplos métodos, fontes e perspectivas antes de
atribuir um score.

\\\\paragraph{Reprodutibilidade e Auditoria}

\\\paragraph{Princípio de Reprodutibilidade Total}

O OpenCode Ecosystem adere ao \textbf{princípio de reprodutibilidade total}:
qualquer afirmação feita nesta dissertação — seja um score, uma correlação,
um gráfico ou uma conclusão — pode ser reproduzida independentemente por um
terceiro com acesso ao repositório GitHub. Este princípio é operacionalizado
através de cinco mecanismos:

\begin{enumerate}[label=(\roman*)]
\item \textbf{Código aberto}: todos os scripts de teste, rastreadores e
      verificadores estão disponíveis em
      \texttt{artigo/evaluations/tests/}.
\item \textbf{Dados abertos}: todos os scores, snapshots e matrizes estão em
      \texttt{cora\_scores.json}, formato JSON human-readable.
\item \textbf{Ambiente reproduzível}: dependências Python são padrão
      (sympy, scipy, numpy) ou implementadas do zero sem dependências externas.
\item \textbf{Trilha de auditoria}: cada score é rastreável a um teste TDD
      específico e a um verificador Cora específico (Anexos A--D).
\item \textbf{Validação externa}: 34 problemas do Project Euler e Rosalind
      fornecem ground truth independente e imutável.
\end{enumerate}

\\\paragraph{Auditoria Independente: Passo a Passo}

Para permitir que uma banca examinadora verifique independentemente qualquer
resultado desta dissertação, documentamos o procedimento completo de auditoria:

\begin{enumerate}[label=\textbf{Passo \arabic*}:,leftmargin=*]
\item \textbf{Clonar o repositório}: \\
      \texttt{git clone https://github.com/MarceloClaro/OpenCode\_Ecosystem}
\item \textbf{Executar o teste exaustivo}: \\
      \texttt{cd artigo/evaluations/tests \&\& python test\_exaustivo\_final.py} \\
      \textit{Esperado: 34/34 PASS, CV=2.2\%}
\item \textbf{Verificar o CORA-Score manualmente}: \\
      \texttt{python cora\_benchmark\_tracker.py --report} \\
      \textit{Esperado: CORA-Score = 3.04, consistente com Anexo D}
\item \textbf{Executar a geometria cognitiva}: \\
      \texttt{python cora\_cognitive\_geometry.py --full} \\
      \textit{Esperado: matriz de dependência, grafo, tensor, curvatura}
\item \textbf{Verificar 5 dimensões em N4}: \\
      \texttt{python cora\_benchmark\_tracker.py --detail D1} (repetir para D2,D3,D7,D10) \\
      \textit{Esperado: cada dimensão com score $\geq$ 3.0}
\end{enumerate}

\\\paragraph{Contraprovas: O Que o Sistema Não Consegue Fazer}

A integridade científica exige que documentemos não apenas os sucessos, mas
também os fracassos — o que o ecossistema \textbf{não} consegue fazer. Esta
seção constitui a ``contraprova'' da dissertação.

\begin{enumerate}[label=(\roman*)]
\item \textbf{EBM com difusão}: o modelo de balanço energético com difusão
      explícita falha por instabilidade numérica (passo de tempo excede o
      critério de estabilidade de Fourier). A versão estacionária (sem
      difusão dependente do tempo) funciona, mas não captura a dinâmica
      transiente. \textbf{Status}: resolvido parcialmente via Crank-Nicolson
      (Seção 9), requer validação adicional.
\item \textbf{Montagem de genoma completo}: o algoritmo de Bruijn implementado
      funciona para genomas sintéticos pequenos ($\sim$3.4kbp), mas não escala
      para genomas bacterianos reais ($\sim$4.6Mbp) sem otimizações de
      memória e estruturas de dados especializadas (Bloom filters, hashing).
      \textbf{Status}: prova de conceito funcional; escala requer engenharia.
\item \textbf{Meta-análise PRISMA}: a extração de metodologias de 50+ artigos
      depende de APIs externas (arXiv, Semantic Scholar) cuja disponibilidade
      e limites de taxa não são controláveis pelo ecossistema. A revisão de
      literatura em escala (D8-N3) permanece o gargalo mais significativo.
      \textbf{Status}: arquitetura definida (chunking + GraphRAG), requer
      implementação e validação com corpus real.
\item \textbf{Simulações HPC}: Schrödinger 2D (split-operator FFT) e
      Navier-Stokes 2D (Re=1000) requerem capacidade computacional além do
      ambiente atual (CPU única, Windows 11). Workarounds via redução
      dimensional existem mas produzem resultados qualitativos, não
      quantitativamente precisos. \textbf{Status}: bloqueio de infraestrutura,
      não de arquitetura.
\end{enumerate}

\\\paragraph{Checklist de Auditoria para a Banca}

\begin{table}[H]\centering\caption{Checklist de verificações para a banca examinadora}
\label{tab:checklist}
\begin{tabular}{p{0.05\textwidth}p{0.55\textwidth}p{0.08\textwidth}p{0.15\textwidth}}
\toprule
\textbf{\#} & \textbf{Verificação} & \textbf{OK?} & \textbf{Evidência} \\
\midrule
1 & CORA-Score = 3.04 (cálculo manual confere) & $\square$ & Anexo D \\
2 & 34/34 testes cegos Project Euler + Rosalind & $\square$ & \texttt{test\_exaustivo\_final.py} \\
3 & Cross-validation K=10 com CV $<$ 5\% & $\square$ & Seção 10.3 \\
4 & 5 dimensões em N4 (D1,D2,D3,D7,D10) & $\square$ & \texttt{cora\_scores.json} \\
5 & Todas as 10 dimensões avaliadas & $\square$ & \texttt{--report} \\
6 & 13 suítes TDD, 113/114 testes GREEN & $\square$ & \texttt{run\_all\_benchmarks.py} \\
7 & Verificadores V1--V7 implementados e testados & $\square$ & Seção 2.4 \\
8 & Matriz de dependência reproduzível ($r$ D1-D10 = 0.97) & $\square$ & \texttt{--matrix} \\
9 & Grafo cognitivo com 4 comunidades & $\square$ & \texttt{--graph} \\
10 & Geodésica de aprendizado: D5$\to$D6$\to$D8$\to$D4$\to$D3 & $\square$ & \texttt{--curvature} \\
\bottomrule
\end{tabular}
\end{table}

\\\\paragraph{Limitações e Agenda de Pesquisa Futura}

\\\paragraph{Limitações Reconhecidas}

Com base nos resultados da validação exaustiva (Seção 10), identificamos
cinco categorias de limitações que afetam a interpretação dos resultados:

\textbf{L1 — Cobertura de domínios}: As 10 dimensões do CORA-Eval cobrem
ciências exatas e da natureza, mas não incluem ciências humanas (economia,
linguística, psicologia) nem engenharias aplicadas (civil, mecânica, elétrica).
A generalização dos resultados para estes domínios não é garantida.

\textbf{L2 — Viés de seleção de tarefas}: As 150 tarefas foram selecionadas
por conveniência (disponibilidade de ground truth), não por amostragem
aleatória estratificada. Isto pode favorecer dimensões com ground truth
abundante (D1, matemática) em detrimento de dimensões com ground truth escasso
(D9, metodologia experimental).

\textbf{L3 — Dependência de verificadores próprios}: Os verificadores V1--V7
foram implementados internamente. Embora a validação externa (Project Euler,
Rosalind) mitigue parcialmente este viés, a auditoria independente dos
verificadores por terceiros fortaleceria a validade dos resultados.

\textbf{L4 — Escala temporal}: Os 8 snapshots cobrem um período de
aproximadamente 11 horas (28--29 de maio de 2026). A extrapolação das
tendências de crescimento (tensor de evolução) para horizontes mais longos
assume estacionariedade que pode não se manter.

\textbf{L5 — Reprodutibilidade entre ambientes}: O ecossistema foi testado
exclusivamente em Windows 11 com Python 3.12. A reprodutibilidade em Linux
ou macOS não foi verificada, embora o código seja portável por design.

\\\paragraph{Agenda de Pesquisa Futura}

\textbf{TF1 —Extensão para ciências humanas}: Adaptar o CORA-Eval para incluir
dimensões de economia (modelagem de equilíbrio geral), linguística
(processamento de corpus) e psicologia experimental (desenho de experimentos
comportamentais).

\textbf{TF2 — Validação externa ampliada}: Identificar fontes de ground truth
externo para física (Olimpíada Internacional de Física), química (Cambridge
Structural Database) e geociências (dados observacionais do IPCC).

\textbf{TF3 — Auditoria independente}: Submeter os verificadores Cora V1--V7
a auditoria por especialistas externos em cada domínio (físicos para V1,
matemáticos para V2, estatísticos para V4, engenheiros de software para V7).

\textbf{TF4 — Escalabilidade computacional}: Migrar simulações HPC (Schrödinger
2D, Navier-Stokes) para ambiente cloud com GPU, e implementar EBM com difusão
implícita (Crank-Nicolson) totalmente validado.

\textbf{TF5 — Meta-análise automatizada}: Implementar o pipeline completo de
revisão sistemática (PRISMA) com integração às APIs arXiv, Semantic Scholar
e OpenAlex, incluindo extração automatizada de effect sizes e funnel plots.

\textbf{TF6 — Geometria diferencial completa}: Implementar conexão de
Levi-Civita, curvatura de Riemann e holonomia no espaço de proficiência,
completando o programa de pesquisa iniciado na Seção 9.

\textbf{TF7 — Estudo longitudinal}: Manter o CORA-Eval ativo por 6--12 meses,
registrando snapshots semanais para validar (ou refutar) as projeções de
crescimento do tensor de evolução e a estabilidade da geodésica de aprendizado.

\\\paragraph{Considerações Éticas e de Transparência}

O OpenCode Ecosystem é um projeto de código aberto (MIT License). Todos os
dados, códigos e resultados apresentados nesta dissertação são públicos e
reproduzíveis. Não há conflitos de interesse a declarar. As referências
bibliográficas incluem DOIs sempre que disponíveis. Os 34 problemas do
Project Euler e Rosalind utilizados para validação externa são de acesso
público e suas respostas são verificáveis independentemente.

A metodologia CORA-Eval não utiliza dados pessoais, não realiza experimentos
com seres humanos ou animais, e não apresenta riscos éticos identificáveis.
O único viés potencial é a sub-representação de pesquisadores do Sul Global
nas fontes de ground truth (Project Euler e Rosalind são plataformas
majoritariamente anglófonas), que deve ser endereçada em trabalhos futuros
através da inclusão de fontes em português, espanhol e outras línguas.
